% ============================================================
% CAPÍTULO 11: RUTAS METABÓLICAS ANABÓLICAS
% ============================================================
% CONTENIDO BASADO EN: Unidad de Aprendizaje III, Tema "Rutas metabólicas anabólicas"
% PROPÓSITO: Diferenciar las rutas anabólicas de las biomoléculas y documentarlas.
% ============================================================

\section{Introducción al anabolismo}

% --- CONSEJO: Define el anabolismo como la fase biosintética del metabolismo, que utiliza energía para construir moléculas complejas a partir de moléculas simples. ---
\resumebox{definicion}{Anabolismo}{Es el conjunto de rutas metabólicas que sintetizan moléculas orgánicas complejas a partir de moléculas más simples, consumiendo ATP y poder reductor (NADPH).}

\subsection{Características generales del anabolismo}
\begin{itemize}
    \item \textbf{Consume energía:} Requiere ATP (de las rutas catabólicas).
    \item \textbf{Utiliza poder reductor:} Requiere NADPH (principalmente de la vía de las pentosas fosfato).
    \item \textbf{Divergencia:} Las rutas anabólicas divergen, utilizando precursores comunes para sintetizar diferentes biomoléculas.
    \item \textbf{Requiere precursores:} Las moléculas simples (precursores) provienen del catabolismo.
\end{itemize}

\section{Biosíntesis de glucógeno}

% --- CONSEJO: El glucógeno es la forma de almacenamiento de glucosa en animales. Se sintetiza a partir de glucosa-6-fosfato y UDP-glucosa. Ocurre principalmente en hígado y músculo. ---
\resumebox{definicion}{Glucogénesis}{Es la ruta anabólica que sintetiza glucógeno a partir de glucosa. Requiere la activación de la glucosa a UDP-glucosa y la formación de enlaces $\alpha$-(1,4) y $\alpha$-(1,6) por la glucógeno sintasa y la enzima ramificadora.}

\section{Biosíntesis de lípidos}

% --- CONSEJO: La síntesis de lípidos (lipogénesis) ocurre principalmente en el hígado y el tejido adiposo. ---
\subsection{Síntesis de ácidos grasos}
% --- CONSEJO: La síntesis de ácidos grasos utiliza acetil-CoA (de la glucosa o de otros precursores) y NADPH. Ocurre en el citosol. La enzima clave es la ácido graso sintasa. ---
\resumebox{definicion}{Lipogénesis}{Es la ruta anabólica que sintetiza ácidos grasos y otros lípidos. Utiliza acetil-CoA como precursor y NADPH como poder reductor.}

\subsection{Síntesis de triglicéridos y fosfolípidos}
% --- CONSEJO: Los ácidos grasos se esterifican con glicerol-3-fosfato para formar triglicéridos (almacenamiento) y fosfolípidos (membranas). ---

\section{Biosíntesis de aminoácidos}

% --- CONSEJO: Los organismos pueden sintetizar todos los aminoácidos que necesitan a partir de intermediarios del metabolismo central (glucólisis, ciclo del ácido cítrico, etc.). Los humanos no pueden sintetizar todos (aminoácidos esenciales). ---
\resumebox{definicion}{Biosíntesis de aminoácidos}{Es el conjunto de rutas anabólicas que sintetizan los 20 aminoácidos a partir de precursores del metabolismo central. Los precursores incluyen intermediarios de la glucólisis, la vía de las pentosas fosfato y el ciclo del ácido cítrico.}

\section{Fotosíntesis, ciclo C3, C4 y fotorrespiración}

% --- CONSEJO: La fotosíntesis es el proceso anabólico por excelencia, que convierte la energía luminosa en energía química (ATP y NADPH) y fija CO2 en carbohidratos. ---
\subsection{Fotosíntesis y ciclo C3 (Ciclo de Calvin)}
% --- CONSEJO: El ciclo de Calvin es la ruta principal de fijación de CO2 en plantas. Utiliza ATP y NADPH de la fase luminosa para fijar CO2 a ribulosa-1,5-bisfosfato (RuBP), formando dos moléculas de 3-fosfoglicerato. ---

\subsection{Ciclo C4 y fotorrespiración}
% --- CONSEJO: El ciclo C4 es una adaptación de algunas plantas para minimizar la fotorrespiración en ambientes cálidos y secos. Concentra CO2 alrededor de la RuBisCO. La fotorrespiración es un proceso ineficiente donde la RuBisCO fija O2 en lugar de CO2. ---

\section{Metanogénesis}

% --- CONSEJO: Es un proceso anabólico y catabólico llevado a cabo por arqueas metanogénicas en ambientes anaeróbicos. Produce metano (CH4) como producto final del metabolismo energético. ---
\resumebox{definicion}{Metanogénesis}{Es un proceso metabólico realizado por arqueas metanogénicas que produce metano (CH4) como producto final de su metabolismo energético. Puede utilizar CO2 y H2, o acetato, como sustratos.}

\section{Regulación del anabolismo}

% --- CONSEJO: El anabolismo está estrechamente regulado para evitar el desperdicio de energía y la acumulación de productos innecesarios. ---
\begin{itemize}
    \item \textbf{Regulación alostérica:} Las enzimas anabólicas son inhibidas por el producto final de la ruta (retroalimentación negativa).
    \item \textbf{Modificación covalente:} La fosforilación de enzimas regula su actividad.
    \item \textbf{Regulación hormonal:} Hormonas como la insulina y el glucagón regulan las rutas anabólicas y catabólicas.
    \item \textbf{Disponibilidad de sustratos:} La célula solo sintetiza las moléculas que necesita en función de la disponibilidad de precursores y la demanda energética.
\end{itemize}

\section{Documentación de las rutas anabólicas}

% --- CONSEJO: Esta sección está diseñada para que el estudiante cree un portafolio de evidencias, diagramando las rutas anabólicas. ---
% --- Se sugiere la creación de mapas metabólicos que integren las rutas anabólicas y catabólicas (por ejemplo, un mapa general del metabolismo). ---